% sections/supervision.tex — Supervision and Causal Time
\section{Supervision and Causal Time}
\label{sec:supervision}

The actor model assumes that any actor may fail, and that the system
must arrange for either restart or destruction of the failed actor
without corrupting the rest of the system.  In a CPU actor system,
supervision is straightforward: the failed actor's mailbox can be
rebuilt from heap, its state can be restored from a checkpoint or
discarded, and its logical clock can continue from wherever it
left off (the OS scheduler is not affected by an application-level
fault).  On a GPU, none of these conveniences hold.  A failed kernel
loses its registers, its shared memory, and its place in the SM
schedule; restoring it means re-launching it from scratch and
re-injecting its persistent state from a snapshot held in HBM.

This section formalizes the supervisor protocol.  We begin with the
Hybrid Logical Clock (HLC) machinery that underwrites causal
ordering on a single device~(\Cref{sec:supervision:hlc}).  We then
describe the four restart policies the supervisor offers
(\Cref{sec:supervision:policies}), give a small-step semantics for
the snapshot--restart cycle (\Cref{sec:supervision:restart}), and
state the main property we prove: HLC monotonicity is preserved
across restart, even though the supervisor must restart from a
snapshot taken at an earlier physical time
(\Cref{thm:hlc-restart-monotone}).  The mechanized counterpart lives
in \texttt{docs/verification/hlc.tla}.

\subsection{Hybrid Logical Clocks}
\label{sec:supervision:hlc}

The HLC, due to Kulkarni et~al.~\citep{kulkarni2014hlc}, combines the
strengths of physical clocks (intuitive wall-clock semantics, useful
for operators) with the strengths of logical clocks
(\citet{lamport1978clocks}; causality preservation under message
passing).  An HLC value is a pair \(\langle p, l\rangle \in \N^2\)
where \(p\) is the physical-time component
(microseconds since epoch, in the runtime; abstracted to a bounded
counter in the spec) and \(l\) is the logical component (a
within-physical-tick counter).

We define two operations.  \(\hlc.\mathsf{tick}\) advances the local
clock when an actor takes a local step:
\[
\hlc.\mathsf{tick}(\langle p, l\rangle, w) \;=\;
  \begin{cases}
    \langle w, 0\rangle             & \text{if } w > p, \\
    \langle p, l + 1\rangle         & \text{otherwise,}
  \end{cases}
\]
where \(w\) is the current wall-clock reading sampled at the moment of
the tick.  \(\hlc.\mathsf{merge}\) advances the local clock on
receipt of a message carrying timestamp \(\langle p_r, l_r\rangle\):
\[
\hlc.\mathsf{merge}(\langle p, l\rangle,
                    \langle p_r, l_r\rangle, w) \;=\;
  \langle p', l'\rangle
\]
where \(p' = \max(p, p_r, w)\) and
\[
  l' \;=\;
  \begin{cases}
    \max(l, l_r) + 1 & \text{if } p' = p = p_r, \\
    l + 1            & \text{if } p' = p \neq p_r, \\
    l_r + 1          & \text{if } p' = p_r \neq p, \\
    0                & \text{otherwise } (p' = w),
  \end{cases}
\]
(matching the four cases in
\texttt{docs/verification/hlc.tla} lines 90--101 and \texttt{HlcClock::update}
in \texttt{crates/ringkernel-core/src/hlc.rs}).

Lexicographic order on HLC values gives a total order on all events
in the system; the standard causality argument
(\citet{kulkarni2014hlc}) shows that if event \(e_1\) happened-before
\(e_2\) in Lamport's sense, then \(e_1\)'s HLC is lexicographically
less than \(e_2\)'s.  We make explicit use of this property in
\Cref{thm:hlc-restart-monotone}.

\begin{property}[Per-Actor Monotonicity]
\label{prop:hlc-per-actor-monotone}
For any actor \(a\) and any two events \(e_1, e_2\) emitted by \(a\)
during a single uninterrupted lifetime (i.e., between any two
adjacent restarts), if \(e_1\) is emitted before \(e_2\), then
\(e_1.\mathit{hlc} <_{\mathrm{lex}} e_2.\mathit{hlc}\).
\end{property}

\begin{proof}
Direct from the definition of \(\hlc.\mathsf{tick}\): the new
physical component is at least the old, and either the physical
strictly increases (and the logical resets to zero, but the pair
strictly increases lexicographically) or the physical is unchanged
and the logical strictly increases.  The argument for
\(\hlc.\mathsf{merge}\) is similar, using \(\max\) on the physical
component.
\qed
\end{proof}

The less immediate question is what happens to monotonicity across
a restart.  After a fault, the supervisor restores the actor from
a snapshot taken at some earlier HLC time \(t_{\mathrm{snap}}\),
and the actor begins emitting events again.  If the new instance
of the actor were to tick its HLC from \(t_{\mathrm{snap}}\), the
events emitted just before the fault (with HLC
\(t_{\mathrm{fault}} > t_{\mathrm{snap}}\)) would be ``in the
future'' relative to the new instance's clock.  This is the
behaviour we need to rule out.

\subsection{Restart Policies}
\label{sec:supervision:policies}

\system{}'s supervisor follows the Erlang/OTP
discipline~\citep{armstrong2007erlang}.  Each supervisor manages a
set of children and applies one of four policies on a child's fault:

\begin{description}
\item[\textsc{Permanent}] Always restart on fault.
\item[\textsc{OneForOne}] Restart only the failed child; leave
siblings running.  This is the default for independent actors.
\item[\textsc{OneForAll}] Restart the failed child and all its
siblings.  Used when siblings share state that the supervisor cannot
prove is consistent after a partial fault.
\item[\textsc{RestForOne}] Restart the failed child and all
siblings created after it (in spawn order).  Used when later
siblings depend on the failed child's invariants.
\end{description}

The choice of policy is a configuration of the supervisor, not of the
child; from the child's point of view, restart is a single semantic
event regardless of which policy provoked it.  The HLC argument below
is therefore independent of which restart policy fired.
\Cref{fig:supervisor-tree} shows a representative supervisor tree
with the four policies in use.

\begin{figure}[htbp]
\centering
\begin{tikzpicture}[
    every node/.style={font=\footnotesize\sffamily},
    sup/.style={draw, rounded corners=2pt, fill=blue!8,
                minimum width=15mm, minimum height=6mm, inner sep=2pt},
    actor/.style={draw, rounded corners=8pt, fill=green!8,
                  minimum width=12mm, minimum height=5mm, inner sep=2pt,
                  font=\scriptsize\sffamily},
    edge/.style={->, >=Stealth, semithick}
  ]
  \node[sup] (root) at (4.5, 2.5) {Root Sup};
  \node[sup] (s1)   at (1.5, 1.0) {Sup A};
  \node[sup] (s2)   at (4.5, 1.0) {Sup B};
  \node[sup] (s3)   at (7.5, 1.0) {Sup C};

  \node[actor] (a1) at (0, -0.5)   {Actor 1};
  \node[actor] (a2) at (1.5, -0.5) {Actor 2};
  \node[actor] (a3) at (3, -0.5)   {Actor 3};

  \node[actor] (b1) at (4.0, -0.5) {Actor 4};
  \node[actor] (b2) at (5.0, -0.5) {Actor 5};

  \node[actor] (c1) at (6.5, -0.5) {Actor 6};
  \node[actor] (c2) at (8.0, -0.5) {Actor 7};
  \node[actor] (c3) at (9.0, -0.5) {Actor 8};

  \draw[edge] (root) -- (s1)
        node[midway, sloped, above, font=\tiny]{OneForOne};
  \draw[edge] (root) -- (s2)
        node[midway, above, font=\tiny]{OneForOne};
  \draw[edge] (root) -- (s3)
        node[midway, sloped, above, font=\tiny]{OneForOne};

  \draw[edge] (s1) -- (a1) node[midway, sloped, above, font=\tiny]{OneForOne};
  \draw[edge] (s1) -- (a2);
  \draw[edge] (s1) -- (a3);

  \draw[edge] (s2) -- (b1) node[midway, sloped, above, font=\tiny]{OneForAll};
  \draw[edge] (s2) -- (b2);

  \draw[edge] (s3) -- (c1) node[midway, sloped, above, font=\tiny]{RestForOne};
  \draw[edge] (s3) -- (c2);
  \draw[edge] (s3) -- (c3);
\end{tikzpicture}
\caption{Representative supervisor tree.  Each supervisor declares a
restart policy that governs how it reacts to a fault among its
children.  \textsc{OneForOne} restarts only the failed child;
\textsc{OneForAll} restarts the failed child and all siblings;
\textsc{RestForOne} restarts the failed child and all siblings
created after it (in spawn order).  A fourth policy,
\textsc{Permanent}, applies per-child rather than per-supervisor.}
\label{fig:supervisor-tree}
\end{figure}

The supervisor also tracks a per-child restart counter
(\(\mathit{rcount}\) in the actor's state vector,
\Cref{def:actor-state}) and a window length.  If the count exceeds a
configured maximum within the window, the supervisor escalates
(forwards the fault to its own supervisor) rather than continuing to
restart.  This bounds the restart-loop blast radius and is the
mechanism by which supervisor trees achieve fault containment.

\subsection{Snapshot--Restart Protocol}
\label{sec:supervision:restart}

A snapshot is taken at every \textsc{Activate} step
(\Cref{rule:activate}): the supervisor copies the actor's
\(\mathit{live}\) bytes into the \(\mathit{snap}\) field of the state
vector.  The snapshot is therefore always equal to the actor's
\(\mathit{live}\) state at the moment of the most recent activation.
Between activations, \(\mathit{snap}\) is immutable; the actor mutates
\(\mathit{live}\) via \textsc{Process} but cannot touch
\(\mathit{snap}\).  This is the snapshot invariant.

On a fault, the supervisor performs three steps:

\begin{enumerate}
\item \textbf{Capture}: the supervisor records the actor's HLC at the
moment of the fault in a per-actor restart state, denoted
\(\mathit{ret}(a) \in \N^2\).  We call this the \emph{retreat
timestamp}.

\item \textbf{Reset}: the actor's \(\mathit{live}\) is overwritten
with \(\mathit{snap}\); \(\mathit{life}\) transitions to
\Initializing{}; \(\mathit{rcount}\) is incremented; the actor's
HLC is set to \(\mathit{ret}(a)\) (i.e., it picks up where it left
off on the physical clock, even though its state has retreated).

\item \textbf{Reactivate}: the actor transitions \Initializing{}
\(\to\) \Active{} via the same \textsc{Activate} rule
(\Cref{rule:activate}), which sets a fresh \(\mathit{snap}\) equal to
the restored \(\mathit{live}\).
\end{enumerate}

The interesting technical move is step~2: although the actor's
\emph{state} retreats to the snapshot's content, the actor's
\emph{HLC} is set to the retreat timestamp \(\mathit{ret}(a)\) (the
HLC value at the moment of the fault), not to whatever value the
HLC had at the moment of the snapshot.  This decouples the data
retreat from the time retreat: the new instance is at the same
logical moment as the failed instance was, with the only difference
being that its computational state is older.

\begin{equation}
\inference[\textsc{Restart}]
{
  \Theta(a) = \langle \Failed{}, \mathit{snap}, \mathit{live}, h, r, \tau\rangle \\
  \mathit{ret}(a) = h_{\mathrm{ret}}
}
{
  \langle \Theta, M, t\rangle
  \trans{\mathsf{restart}(a)}
  \langle \Theta', M, t+1\rangle
}
\label{rule:restart}
\end{equation}

\noindent where
\(\Theta' = \Theta[a \mapsto
  \langle \Initializing{}, \mathit{snap}, \mathit{snap},
          h_{\mathrm{ret}}, r+1, \tau\rangle]\).

\subsection{HLC Monotonicity Across Restart}
\label{sec:supervision:monotonicity}

The main property of the supervisor protocol is that HLC
monotonicity—\Cref{prop:hlc-per-actor-monotone}—survives restart.
Concretely: there is no \emph{observable} HLC value emitted by a
restarted instance that is less-than-or-equal to an HLC value emitted
by the failed instance.  Informally: \emph{restart may roll back
state, but it must not roll back time.}  The retreat-timestamp
protocol of \Cref{sec:supervision:restart} is exactly the mechanism
that keeps those two rollbacks separate.

\begin{theorem}[HLC Restart Monotonicity]
\label{thm:hlc-restart-monotone}
Let \(a\) be an actor that fails at HLC time
\(h_{\mathrm{fault}} \in \N^2\) and is restarted by
\textsc{Restart}~(\Cref{rule:restart}).  For every event \(e\)
emitted by the restarted instance of \(a\) (i.e., by any
\textsc{Process} step taken from a configuration \(C'\) reachable
after the \textsc{Restart} transition), we have
\(h_{\mathrm{fault}} <_{\mathrm{lex}} e.\mathit{hlc}\).
\end{theorem}

\begin{proof}[Proof sketch]
Let \(h_{\mathrm{ret}}\) be the retreat timestamp captured by the
supervisor at the moment of fault; by step~1 of the protocol,
\(h_{\mathrm{ret}}\) is the HLC value the failed instance held at
the moment of fault, which is the HLC value of the most recent
event emitted by the failed instance.  Hence \(h_{\mathrm{fault}}
\le_{\mathrm{lex}} h_{\mathrm{ret}}\).  By step~2 of the
protocol~(\Cref{rule:restart}), the restarted instance's HLC is
initialized to \(h_{\mathrm{ret}}\).  Any subsequent event \(e\)
emitted by the restarted instance is the result of either a
\(\hlc.\mathsf{tick}\) or a \(\hlc.\mathsf{merge}\) from a clock
value that starts at \(h_{\mathrm{ret}}\); by
\Cref{prop:hlc-per-actor-monotone} applied to the new instance's
lifetime, every such event has HLC strictly greater than
\(h_{\mathrm{ret}}\).  Therefore \(e.\mathit{hlc}
>_{\mathrm{lex}} h_{\mathrm{ret}} \ge_{\mathrm{lex}}
h_{\mathrm{fault}}\), as required.
\qed
\end{proof}

The hypothesis is that the supervisor faithfully captures
\(\mathit{ret}(a)\) before allowing the restart to proceed; if the
supervisor were to skip step~1 and initialize the new instance at
\(\hlc.\mathsf{tick}(h_{\mathrm{snap}}, w)\) (i.e., from the
snapshot's HLC, for any wall-clock reading \(w\) at the moment of
restart), the property would fail in the regime where
\(h_{\mathrm{snap}} <_{\mathrm{lex}} h_{\mathrm{fault}}\) and
\(w\) does not exceed the physical component of
\(h_{\mathrm{fault}}\): a single tick is then not enough to
overshoot \(h_{\mathrm{fault}}\) when the physical components
match.  The runtime enforces step~1
explicitly: the supervisor's restart logic refuses to issue
\textsc{Restart} until \(\mathit{ret}(a)\) has been written to the
durable supervisor state.

\subsection{Causal Ordering Across Restart}
\label{sec:supervision:causal}

\Cref{thm:hlc-restart-monotone} has an immediate corollary in terms
of message-passing causality:

\begin{corollary}[Restart-Safe Causal Ordering]
\label{cor:restart-causal}
Let \(e_1\) be a message emitted by the failed instance of \(a\) and
\(e_2\) be a message emitted by the restarted instance of \(a\).
For any actor \(b\) that receives both \(e_1\) and \(e_2\), the HLC
order \(e_1.\mathit{hlc} <_{\mathrm{lex}} e_2.\mathit{hlc}\) holds,
and hence \(b\)'s
\(\hlc.\mathsf{merge}\) operation will assign \(e_2\) a strictly
greater HLC than \(e_1\).
\end{corollary}

\begin{proof}
Direct from \Cref{thm:hlc-restart-monotone} by transitivity:
\(e_1.\mathit{hlc} \le_{\mathrm{lex}} h_{\mathrm{fault}}
<_{\mathrm{lex}} e_2.\mathit{hlc}\), where the first inequality
holds because \(e_1\) was emitted by the failed instance no later
than the moment of fault.  The merge property follows from the
monotonic-merge property of the HLC
(\citet{kulkarni2014hlc}, Theorem~3).
\qed
\end{proof}

\Cref{cor:restart-causal} is the property that lets external
observers (e.g., the audit trail of \Cref{sec:tenancy:audit}) treat
a restarted actor's messages as causally subsequent to the failed
instance's messages, even when the actor's data has retreated.  This
is the analogue, in our setting, of the durability guarantees of
write-ahead logging in a CPU database.

\subsection{Mechanization}
\label{sec:supervision:mechanization}

The HLC machinery is mechanized in \TLA{} in
\texttt{docs/verification/hlc.tla}.  The mechanization keeps the
HLC values, the wall-clock counter, and a sequence of emitted
events; it model-checks the invariants \texttt{PhysicalTimeMonotonic},
\texttt{CausalOrdering}, \texttt{PerNodeMonotonic}, and
\texttt{EventCountsBounded}.  The restart rule of
\Cref{rule:restart} is not represented directly in
\texttt{hlc.tla}—the spec is restricted to single-instance per node
HLC dynamics—but the property of \Cref{thm:hlc-restart-monotone}
follows from \texttt{PerNodeMonotonic} and the protocol invariant
that \(\mathit{ret}(a)\) is captured before \textsc{Restart},
provable analytically as above.  The actor-lifecycle spec in
\texttt{actor\_lifecycle.tla} models the snapshot--restart cycle
without explicit HLC; the two specs together cover the property as
stated.  Folding the two specs together so that TLC can check
\Cref{thm:hlc-restart-monotone} directly is a target for v1.1 of
the verification suite (see \Cref{sec:discussion}).
